% SPDX-License-Identifier: CC-BY-4.0
% P4 coupled frame+timescale gauge manuscript. Target: NAVIGATION (ION) / Journal of Geodesy / Metrologia.
% Build: tectonic main.tex   (or pdflatex + bibtex + pdflatex x2)
\documentclass[11pt]{article}

\usepackage[utf8]{inputenc}
\usepackage[T1]{fontenc}
\usepackage{lmodern}
\usepackage[margin=1in]{geometry}
\setlength{\emergencystretch}{3em}
\usepackage{amsmath}
\usepackage{amssymb}
\usepackage{amsthm}
\usepackage{bm}
\usepackage{graphicx}
\usepackage{booktabs}
\usepackage{array}
\usepackage{siunitx}
\usepackage{xcolor}
\usepackage[colorlinks=true,linkcolor=black,citecolor=black,urlcolor=blue]{hyperref}
\usepackage[numbers,sort&compress]{natbib}
\usepackage{orcidlink}

\graphicspath{{./}{figures/}}
\setlength{\abovecaptionskip}{9pt}
\setlength{\belowcaptionskip}{6pt}

% Theorem environments
\newtheorem{theorem}{Theorem}
\newtheorem{lemma}{Lemma}
\newtheorem{proposition}{Proposition}
\newtheorem{corollary}{Corollary}
\theoremstyle{definition}
\newtheorem{definition}{Definition}
\newtheorem{remark}{Remark}

% Notation macros
\newcommand{\code}[1]{\texttt{#1}}
\newcommand{\R}{\mathbb{R}}
\newcommand{\Imat}{\bm{I}}                       % Fisher information matrix
\newcommand{\Gmat}{\bm{G}}                       % geometry (design) matrix
\newcommand{\thetav}{\bm{\theta}}                % coupled datum vector
\newcommand{\Umat}{\bm{U}}                       % null-space basis (9 x d)
\newcommand{\Nsp}{\mathcal{N}}                   % null space
\newcommand{\Ssp}{\mathcal{S}}                   % spatial subspace
\newcommand{\Tsp}{\mathcal{T}}                   % temporal subspace
\newcommand{\Pnull}{\bm{P}_{\!\Nsp}}            % null-space projector U U^T
\newcommand{\Pst}{\bm{P}_{\!\Ssp\Tsp}}          % spatial-temporal off-diagonal block
\newcommand{\dtau}{\delta\tau}                   % timescale offset gauge
\newcommand{\dalpha}{\delta\alpha}               % timescale rate gauge
\newcommand{\rmoon}{R_{\leftmoon}}               % lunar reference radius
\newcommand{\Umoon}{U_{\leftmoon}}               % lunar surface potential
\newcommand{\gmoon}{g_{\leftmoon}}               % lunar surface gravity
\newcommand{\rng}{\operatorname{range}}
\newcommand{\Sig}{\bm{\Sigma}}                   % covariance
\DeclareMathOperator{\rank}{rank}

% Honesty tags for proven / validated / representative-modelled quantities
\newcommand{\Proven}{\textsc{Proven}}
\newcommand{\Modelled}{\textsc{Modelled}}
\newcommand{\Validated}{\textsc{Validated}}
\newcommand{\modtag}{{\footnotesize\textsf{[M]}}}

\title{\textbf{The coupled frame--timescale gauge of lunar PNT: a
basis-invariant datum-defect classification, the two frame--time couplings,
and why a lunar reference \emph{system} must fix both a spatial datum and a
temporal convention}}

\author{Chakshu Baweja\,\orcidlink{0009-0008-2098-0751}\\
\small Ashforde O\"U, Tallinn, Estonia \\
\small \texttt{contact@ashforde.org}}

\date{\today}

\begin{document}
\maketitle

\begin{abstract}
A lunar positioning, navigation, and timing (PNT) realization jointly fixes a spatial
reference \emph{frame} and a \emph{timescale}, and the two are not independent. We
classify the null space---the datum defect---of the \emph{coupled} frame$\oplus$timescale
Fisher information of a lunar network over a nine-parameter datum (three translations, a
scale, three rotations, a timescale offset $\dtau$, and a timescale rate $\dalpha$), and
we separate throughout what we cite, prove, validate, and model. Under the SI second and
the $c$-defined metre---adopted up front to quarantine the definitional
unit-of-time/scale coupling and to keep scale \emph{observable} rather than a free datum
direction---we prove (Theorem~1) that the fully internal \emph{free-network} datum gauge of
a one-way lunar ranging problem---mutual common-mode clocks, no external tie---is exactly
\emph{eight}: three translations, three rotations, $\dtau$, and $\dalpha$; the tied,
anchored realization we validate has defect zero. The classification is \Proven{}
\emph{basis-invariant}---a function of the null \emph{subspace}
through the projector $\Pnull=\Umat\Umat^\top$, not of any rotation-ambiguous null
basis---and a coupled defect is certified only by the off-diagonal projector block
$\Pst\neq\bm0$: a non-block-diagonal Fisher is necessary but not sufficient. We quantify
the two frame--time couplings (\Modelled{} magnitudes). Coupling~1,
range$\leftrightarrow$clock-offset, has unit tie
$1\,\mathrm{m}\leftrightarrow3.336\,\mathrm{ns}$ ($=1/c$); we correct a natural but false
claim by showing that in a geometrically diverse multi-beacon one-way network the Fisher
is full-rank and two-way ranging gives \emph{no} selective scale$\leftrightarrow$offset
decoupling beyond data volume (lift ratio $\approx1.00$, computed on the modelled network).
Coupling~2,
relativistic rate$\leftrightarrow$frame, is fundamentally weak: the frame-estimation datum
perturbs the lunar timescale rate by $\leq\!\sim\!2\times10^{-17}$, against a
$\pm1.7\times10^{-11}$ reference-\emph{surface} definitional band---the rate is
gauge-\emph{robust} to estimation but gauge-\emph{sensitive} to definition. Our single
\Validated{} claim is that the coupled-Fisher linear-algebra pipeline reproduces an
independent numpy/LAPACK computation, to relative and absolute error below $10^{-3}$, on
observation rows carrying real DE440 lunar-orientation partials, for a defect-$0$ and an
exact-rate-defect network. Two design laws follow: geometric diversity separates the
frame scale from the timescale offset; and the rate's residual gauge is definitional---so
a lunar reference \emph{system} standard must fix both a spatial datum and a temporal
convention.

\end{abstract}

\noindent\textbf{Keywords:} lunar PNT; geodetic datum; gauge freedom; null space;
Fisher information; reference frame; timescale; relativistic clock rate; LunaNet;
Moonlight; LCNS; DE440; identifiability

\vspace{0.5em}
% The \rmoon/\Umoon/\gmoon macros in main.tex expand to R_{\leftmoon} etc.; the
% \leftmoon glyph needs a symbol package not loaded in the (fixed) preamble. Redefine
% here (document body, before first use in section 2) with a plain "M" (Moon) subscript.
\renewcommand{\rmoon}{R_{\mathrm{M}}}
\renewcommand{\Umoon}{U_{\mathrm{M}}}
\renewcommand{\gmoon}{g_{\mathrm{M}}}

\section{Introduction}
\label{sec:intro}

The next decade will place several independent positioning, navigation, and timing (PNT)
services around the Moon---NASA's LunaNet and its Lunar Communications Relay and
Navigation Systems, ESA's Moonlight programme and its Lunar Communications and Navigation
Service (LCNS), and their national counterparts \citep{israel2020lunanet,%
fienga2024moonlight,lnis2023}. Each such service realizes not one reference object but
\emph{two}: a spatial reference \emph{frame} (an origin, an orientation, and a scale for
lunar coordinates) and a \emph{timescale} (a lunar coordinate time, its offset from
Terrestrial Time, and its rate). The two are realized by the \emph{same} measurements---%
one-way and two-way ranges between orbiters and surface beacons carry both the geometry
and the clocks---so their datum freedoms are, in principle, coupled. The design of a
lunar reference \emph{system}, now under active study by the IAU and in the Moonlight
lunar-reference-frame and lunar-time work \citep{fienga2024moonlight}, must therefore ask
a joint question: what does a lunar PNT network \emph{fail} to determine about the
combined frame and timescale, and which of those failures are spatial, which temporal,
and which genuinely \emph{coupled}?

This paper answers that question by classifying the null space---the geodetic
\emph{datum defect}, or gauge freedom \citep{baarda1973,teunissen2000}---of the
\emph{coupled} frame$\oplus$timescale Fisher information, over a nine-parameter datum that
carries the spatial Helmert seven of our companion datum-identifiability paper
\citep{baweja2026datum} together with two temporal parameters, a clock offset $\dtau$ and
a clock rate $\dalpha$. It is the capstone of a companion arc: P1 classifies the
\emph{spatial} datum \citep{baweja2026datum}; P2 the cross-provider interoperability floor
\citep{baweja2026interop}; P3 the common-mode error class an integrity monitor cannot see
\citep{baweja2026commonmode}; and this paper (P4) adds \emph{time} and, crucially, the
\emph{coupling}. We are deliberate, throughout, about the status of every statement---%
cited, proven, validated, or modelled---because the naive versions of two of the headline
claims are false, and a careful reviewer will (rightly) reject them.

\paragraph{The SI-second convention, adopted up front.}
We fix the unit convention at the outset: the SI second (Cs) and the $c$-defined metre.
This choice is not cosmetic; it \emph{quarantines a definitional trap}. A pseudorange is
$c\,\Delta t$, so a drift of the \emph{unit of time} $\delta\alpha$ masquerades as a
fractional scale $\delta s$ at coefficient $\approx1$: the two are indistinguishable
\emph{if the unit of time is left free}. This is a choice of units, not an estimation
degeneracy, and adopting the SI second removes it once and for all. Under SI, a
pseudorange is scale-bearing through the fixed $c$, so \textbf{scale is observable}---it
is \emph{not} a datum defect, and we never treat it as a free datum direction. The residual
temporal
freedom $\dalpha$ is then the \emph{estimation-deviation} freedom about the
physically-determined secular rate of lunar coordinate time relative to Terrestrial Time,
not the physical rate itself; we word this carefully wherever it appears.

\paragraph{What we prove.}
The core contribution is a classification theorem for the coupled datum defect
(Section~\ref{sec:theorem}), stated in a \emph{basis-invariant} form. The datum defect is
a property of the null \emph{subspace} of the Fisher information, and we prove that its
spatial/temporal/coupled decomposition is a function of that subspace through the
orthogonal projector $\Pnull=\Umat\Umat^\top$, hence invariant under the
rotation-ambiguous choice of null basis $\Umat\mapsto\Umat\bm Q$. This invariance is the
rigor upgrade of the paper: the tempting per-vector measure ``energy of a null vector in
the temporal block'' is \emph{ill-posed}, because a null basis can be rotated freely. We
then prove (Theorem~\ref{thm:classify}) that a fully internal \emph{free-network} one-way
lunar ranging problem under SI units---mutual common-mode clocks, no external tie---has an
exact datum gauge of \emph{eight}: three translations, three rotations, the clock offset
$\dtau$, and the clock rate $\dalpha$, with scale observable. The joint gauge is therefore
\emph{not} the direct sum of the spatial and temporal gauges; and external ties lift each
temporal defect (Theorem~\ref{thm:lift}), carrying the free-network gauge of eight down to
the defect-$0$ tied realization we validate (Section~\ref{sec:anchor}). A sharp corollary bounds what a Fisher's off-diagonal
block can and cannot certify: a non-block-diagonal Fisher is \emph{necessary but not
sufficient} for a coupled \emph{defect} (Remark~\ref{rem:necessary}).

\paragraph{What we model, and the two couplings.}
The network geometry (beacons and orbits) is \Modelled{}---representative, not flown---and
so are the two coupling \emph{magnitudes} (Section~\ref{sec:couplings}), which are
first-principles engineering bounds cross-checked against the published lunar-clock-rate
band, not certified budgets. \emph{Coupling~1} (range$\leftrightarrow$clock-offset) is
operationally strong but a \emph{near}-defect, not an exact one; its classical root is the
GNSS radial/clock degeneracy, which we cite and do not claim
\citep{sturza1988,brown1992,parkinson1988}. Here we state, and correct, a claim whose
naive form is false: it is \emph{not} true that two-way ranging generically decouples the
lunar frame scale from the timescale offset at the datum level. In a geometrically diverse
multi-beacon multi-epoch one-way network the coupled Fisher is full-rank, the
$\{s,\dtau\}$ marginal is positive-definite, and adding two-way ranges lifts the marginal
by the same factor as adding an equal count of \emph{more} one-way ranges (lift ratio
$\approx1.00$, computed on the modelled network); the strong
$1\,\mathrm{m}\leftrightarrow3.336\,\mathrm{ns}$ degeneracy is a
per-user, poor-local-geometry effect that does not aggregate into a datum-level
near-defect. \emph{Coupling~2} (relativistic rate$\leftrightarrow$frame) is fundamentally
weak: the frame estimation datum perturbs the lunar timescale rate by at most
$\sim\!2\times10^{-17}$, some six orders of magnitude below the
$\pm1.7\times10^{-11}$ reference-\emph{surface} definitional band. The rate is
gauge-\emph{robust} to estimation and gauge-\emph{sensitive} to definition.

\paragraph{What we validate on real data.}
Exactly one claim is externally \Validated{} (Section~\ref{sec:anchor}): that the
coupled-Fisher linear-algebra pipeline---preconditioned assembly, eigendecomposition,
basis-invariant null-space classification, and the $\{s,\dtau\}$ Schur marginal---%
reproduces an \emph{independent} numpy/LAPACK computation to relative and absolute error
below $10^{-3}$, on observation rows whose spatial partials carry the \emph{real} DE440
lunar physical-axis orientation \citep{park2021de440}. The two anchored networks are a
defect-$0$ full-rank network and an exact rate-defect network; both have zero coupled
dimension, so the nonzero-coupling ($\Pst\neq\bm0$) branch of the classifier is
\Proven{} by construction and covered by module unit tests---including a rotated-basis
invariance test---and is \emph{not} claimed as externally validated. We state this scope
precisely, because it is the discipline the paper is built on.

\paragraph{Honest-novelty statement.}
To fix the boundary once: the geodetic datum defect and its null-space/gauge machinery are
classical \citep{baarda1973,teunissen2000}; the range$\leftrightarrow$clock (radial/clock)
degeneracy is classical GNSS \citep{sturza1988,brown1992,parkinson1988}; and the
relativistic dependence of a clock's rate on the gravitational potential it sits in is
standard relativistic geodesy and the IAU relativistic timescale framework
\citep{soffel2003iau,mueller2018clocks}. We claim novelty for none of these. Our
contributions are: (i) the \emph{coupled} frame$\oplus$timescale datum-defect
classification in basis-invariant form, with the exact free-network defect count of eight
under SI and the necessary-but-not-sufficient off-diagonal criterion; (ii) the corrected,
datum-level finding that two-way ranging provides no selective scale$\leftrightarrow$offset
decoupling in a diverse network; (iii) the quantified rate$\leftrightarrow$frame coupling
bound ($\leq\!\sim\!2\times10^{-17}$) and the gauge-robust-to-estimation /
gauge-sensitive-to-definition dichotomy; and (iv) the resulting two design laws for a
lunar reference-system standard.

\paragraph{Honesty scope.}
We separate four kinds of statement and label them at every occurrence. The datum-defect
and gauge-freedom machinery, the radial/clock degeneracy, and the relativistic clock-rate
physics are \emph{cited} prior art. The classification theorem, its basis-invariance, the
exact defect count, the necessary-not-sufficient criterion, and the lift results
(Section~\ref{sec:theorem}) are \Proven{} algebra. The single coupled-Fisher
linear-algebra reproduction on real DE440-orientation rows (Section~\ref{sec:anchor}) is
\Validated{} against an independent numerical oracle. Everything else---the network
geometry and the two coupling \emph{magnitudes} (Section~\ref{sec:couplings})---is
\Modelled{} and tagged \modtag{}. There is no flight data, no certified timescale budget,
no technology-readiness claim above analysis, and no single-number operational rate claim;
every number is regenerated from the open engine by one command
(Section~\ref{sec:availability}).

\section{The coupled datum, the observation model, and the coupled Fisher information}
\label{sec:model}

We work with a linearized model in which a lunar PNT network of orbiters and surface
beacons estimates, jointly, a spatial reference frame and a timescale. This section fixes
the unit convention, defines the nine-parameter coupled datum, gives the four
observation-row builders, and assembles the coupled Fisher information whose null space
Section~\ref{sec:theorem} classifies. All quantities named \code{like\_this} are the
corresponding items in the open engine module \code{lunar\_gauge}
(Section~\ref{sec:availability}).

\subsection{The SI-second convention}
\label{sec:si}

Throughout we adopt the SI second (realized by the caesium hyperfine transition) and the
metre defined by the fixed speed of light $c$. A pseudorange is the product $c\,\Delta t$
of this fixed $c$ and a measured time interval. The consequence is decisive for the datum
analysis: because $c$ is fixed, a pseudorange \emph{carries} scale, and a fractional scale
change $\delta s$ of the lunar body produces an observable change in every range. Scale is
therefore \textbf{observable} and is \emph{not} a datum defect.

Had we instead left the unit of time free, a drift $\delta\alpha$ of that unit would
rescale every measured interval and hence every range, entering the observations
identically to a fractional scale $\delta s$ at coefficient $\approx1$; the two would be
indistinguishable. That degeneracy is a matter of \emph{units}, not of estimation
geometry, and adopting the SI second removes it once and for all. We therefore keep two
distinct temporal parameters below---a clock \emph{offset} $\dtau$ and a clock
\emph{rate} $\dalpha$---with $\dalpha$ understood as the estimation-deviation freedom
about the physically-determined secular rate of lunar coordinate time relative to
Terrestrial Time (TT), \emph{not} the physical rate itself. We treat scale as an
observable throughout, never as a free datum direction.

\subsection{The nine-parameter coupled datum}
\label{sec:datum}

The coupled datum vector is the nine-component
\begin{equation}
  \thetav
  \;=\;
  \bigl[\,\underbrace{t_x,\;t_y,\;t_z}_{\text{translation }\bm t},\;
         \underbrace{s}_{\text{scale}},\;
         \underbrace{\theta_x,\;\theta_y,\;\theta_z}_{\text{rotation }\bm\theta},\;
         \underbrace{\dtau}_{\text{offset}},\;
         \underbrace{\dalpha}_{\text{rate}}\,\bigr]^\top
  \;\in\;\R^{9},
  \label{eq:datum9}
\end{equation}
indexed $0,\dots,8$. Indices $0$--$2$ are a translation $\bm t$ of the lunar frame origin
(metres); index $3$ is the fractional scale $s$; indices $4$--$6$ are an infinitesimal
rotation $\bm\theta$ of the frame (radians); index $7$ is the timescale offset $\dtau$;
and index $8$ is the timescale rate $\dalpha$. In the engine these positions are the
constants $\code{IDX\_SCALE}=3$, $\code{IDX\_OFFSET}=7$, $\code{IDX\_RATE}=8$, with
$\code{N\_GAUGE}=9$. The first seven components---translation, scale, and rotation---form
the \emph{spatial} block; concretely the spatial subspace and the temporal subspace are
\begin{equation}
  \Ssp \;=\; \operatorname{span}\{\bm e_0,\dots,\bm e_6\},
  \qquad
  \Tsp \;=\; \operatorname{span}\{\bm e_7,\bm e_8\},
  \label{eq:subspaces}
\end{equation}
with $\bm e_k$ the $k$-th standard basis vector of $\R^{9}$; $\Ssp$ carries the Helmert
seven-parameter spatial datum of \citep{baweja2026datum} (three translations, three
rotations, and the now-observable scale), and $\Tsp$ carries the two temporal parameters.
The clock offset $\dtau$ sits at the boundary index $7$: it is a temporal parameter, but
because a common node clock offset enters a one-way range exactly as a common range bias
it is the natural coupling handle between $\Ssp$ and $\Tsp$, as Section~\ref{sec:couplings}
makes precise.

\subsection{Two network models: the free internal gauge and the anchored realization}
\label{sec:models}

The datum \eqref{eq:datum9} is classified in two operationally distinct settings, and the
distinction is decisive for the defect count; we fix it here once.

\paragraph{(A) The free (internal) network---the datum gauge.}
In the \emph{free} setting the network's own nodes carry mutual, \emph{common-mode} clocks
and no node position or time is known externally. A one-way pseudorange then carries the
clock term as a \emph{difference} (receiver minus transmitter), so a common network
timescale offset $\dtau$ and a common rate $\dalpha$ cancel identically in every internal
observation: the two datum columns vanish in the differential design,
\begin{equation}
  g^{\mathrm{free}}_{\dtau}\;=\;0,\qquad g^{\mathrm{free}}_{\dalpha}\;=\;0
  \qquad\text{(common mode)},
  \label{eq:free}
\end{equation}
so $\bm e_{\dtau}=\bm e_7$ and $\bm e_{\dalpha}=\bm e_8$ lie in $\ker\Imat$ exactly. This
free-network design---the temporal analogue of the P1 spatial free-network datum defect
\citep{baweja2026datum}, extended to time---defines the coupled datum \emph{gauge} that
Theorem~\ref{thm:classify} classifies. It is not the design used for the validated anchor.

\paragraph{(B) The anchored network---the realized tie.}
In the \emph{anchored} setting the orbiters sit at \emph{known} inertial ephemeris
positions with a \emph{known} time reference; this is the model of the one-way row
\eqref{eq:oneway} used for the validated anchor of Section~\ref{sec:anchor}. Here the
offset column is $+1$---an \emph{absolute} offset measured against that known reference,
not a differenced quantity---and the rate column is the absolute accumulation
$t_{\mathrm e}/T_{\mathrm B}$. The known orbiter frame and time reference constitute an
implicit external frame/time tie, so the timescale offset and (with a nonzero time
baseline) the rate become \emph{observable}: a geometrically diverse anchored network
observes the whole datum, with defect $0$ (Section~\ref{sec:anchor}). The free-network
gauge of eight (Theorem~\ref{thm:classify}) and the tied realization of defect $0$ are
connected by the lift results of Theorem~\ref{thm:lift}.

\subsection{The four observation-row builders}
\label{sec:rows}

Each scalar observation contributes one row $\bm g\in\R^{9}$, the gradient of the modelled
observable with respect to $\thetav$. There are four row types.

\paragraph{One-way range (\code{oneway\_range\_row}).}
For a one-way orbiter$\to$beacon range accumulated over an elapsed time $t_{\mathrm e}$,
\begin{equation}
  \bm g^{\mathrm{ow}}
  \;=\;
  \bigl[\,\underbrace{d_0,\dots,d_6}_{\text{spatial datum-7 partials}},\;
          \underbrace{1}_{\dtau},\;
          \underbrace{t_{\mathrm e}/T_{\mathrm B}}_{\dalpha}\,\bigr],
  \label{eq:oneway}
\end{equation}
where the spatial columns $0$--$6$ are the line-of-sight-projected Helmert
datum-7 partials of \citep{baweja2026datum}, computed through the \emph{real} DE440
lunar physical-axis (PA) orientation path (the engine routine
\code{orbiter\_range\_row\_datum7}); the offset column is unity because, in this
\emph{anchored} form (model (B) of Section~\ref{sec:models}), the clock offset $c\,\dtau$
is measured \emph{absolutely} against the known orbiter time reference and enters the range
directly---in contrast to the free/differential form \eqref{eq:free}, where a common offset
cancels; and the rate column is
$t_{\mathrm e}/T_{\mathrm B}$, since a rate offset accumulates linearly with elapsed time,
divided by a baseline $T_{\mathrm B}=\SI{86400}{s}$ (one Julian day) purely to bring the
column to $\mathcal{O}(1)$ (Section~\ref{sec:fisher}).

\paragraph{Two-way range (\code{twoway\_range\_row}).}
A two-way (round-trip) range shares the spatial columns of \eqref{eq:oneway} at the same
geometry, but its temporal columns are \emph{both} zero,
\begin{equation}
  \bm g^{\mathrm{tw}}
  \;=\;
  \bigl[\,d_0,\dots,d_6,\;0,\;0\,\bigr],
  \label{eq:twoway}
\end{equation}
because the round trip cancels both the common node clock offset and its rate. Two-way
ranging thus contributes spatial information without touching $\dtau$ or $\dalpha$; whether
this selectively lifts the offset$\leftrightarrow$rate freedom is the subject of the
corrected finding in Section~\ref{sec:coupling1}.

\paragraph{External time tie (\code{time\_tie\_row}).}
An external common-view Earth--Moon time comparison observes the offset directly,
\begin{equation}
  \bm g^{\mathrm{tt}}
  \;=\;
  \bigl[\,0,0,0,0,0,0,0,\;1,\;0\,\bigr]
  \;=\; \bm e_7 ,
  \label{eq:timetie}
\end{equation}
lifting the pure temporal-offset defect.

\paragraph{Relativistic rate tie (\code{rate\_tie\_row}).}
The physical statement ``the realized clock rate is a function of the frame scale'' is the
rank-one tie
\begin{equation}
  \bm g^{\mathrm{rt}}
  \;=\;
  \bigl[\,0,0,0,\;-\tfrac{\partial\alpha}{\partial s},\;0,0,0,\;0,\;1\,\bigr],
  \label{eq:ratetie}
\end{equation}
with a unit rate column and a scale column $-\partial\alpha/\partial s$ set by the
relativistic Jacobian of Section~\ref{sec:coupling2}. It converts the rate defect into a
coupled scale$\leftrightarrow$rate defect \emph{unless} scale is independently tied
(Theorem~\ref{thm:lift}). The velocity and radial couplings of Section~\ref{sec:coupling2}
are real relativistic effects but do not map onto the nine-vector datum---there is no
velocity gauge, and a uniform frame translation cancels on a symmetric shell---so they
enter the physical rate budget, not a datum-level tie row.

\subsection{The coupled Fisher information and its preconditioning}
\label{sec:fisher}

With independent Gaussian observations of variance $\sigma^2$, the coupled Fisher
information over $\thetav$ is the $9\times9$ Gram sum
\begin{equation}
  \Imat(\thetav)
  \;=\;
  \sum_i \frac{1}{\sigma_i^2}\,\bm g_i\,\bm g_i^\top
  \;\in\;\R^{9\times9},
  \label{eq:fisher}
\end{equation}
assembled block-by-block by \code{assemble\_coupled\_info}. Because the raw partials span
many orders of magnitude---translation columns are $\mathcal{O}(1)$ line-of-sight
components, while scale and rotation columns carry a lever arm of order the lunar radius---%
the assembler \emph{preconditions} by dividing the scale and three rotation columns
(indices $3$--$6$) by $\rmoon=\SI{1737400}{m}$, bringing all nine partials to
$\mathcal{O}(1)$ and reducing the condition number. Preconditioning is a positive diagonal
column scaling $\Imat\mapsto\bm D\Imat\bm D$ with $\bm D\succ0$; it changes eigenvalues but
\emph{not} the null-space structure, which is the object we classify. We record this
invariance and use it below.

\begin{lemma}[Preconditioning preserves the datum-defect counts]
\label{lem:precond}
Let $\bm D=\operatorname{diag}(d_0,\dots,d_8)$ with every $d_k>0$ act within $\Ssp$ and
within $\Tsp$ separately (as the preconditioner does: it scales only indices $3$--$6$, all
in $\Ssp$). Then $\dim\ker(\bm D\Imat\bm D)=\dim\ker\Imat$, and $\bm v\mapsto\bm D^{-1}\bm v$
is a bijection $\ker(\bm D\Imat\bm D)\!\leftrightarrow\!\ker\Imat$. Consequently the four
\emph{integer} counts of the classification of Section~\ref{sec:theorem}---the defect $d$,
the two subspace dimensions, and $\operatorname{coupled\,dim}$---are unchanged, as is the
binary \emph{zeroness} certificate $\Pst=\bm0$ (block-diagonality of the null projector).
The projector-block \emph{norm value} $\lVert\Pst\rVert_{\mathrm F}$ is \emph{not} preserved:
$\bm D$ is a non-orthogonal change of basis, so that value is invariant only under the
orthogonal basis changes of Proposition~\ref{prop:invariance}, not under the preconditioner.
\end{lemma}

\begin{proof}
$\bm D\Imat\bm D\,\bm w=\bm 0 \iff \Imat(\bm D\bm w)=\bm 0$ since $\bm D\succ0$ is
invertible, so $\bm w\in\ker(\bm D\Imat\bm D)\iff\bm D\bm w\in\ker\Imat$; the map is a
linear bijection, giving equal dimensions. The preconditioner scales indices $3$--$6$,
all in $\Ssp$, so it preserves the split \eqref{eq:subspaces} and commutes with the
spatial/temporal projections used by the classifier; the rank formulas of
Section~\ref{sec:theorem}, being functions of the images of $\ker$ under the fixed
projections onto $\Ssp$ and $\Tsp$, are therefore unchanged. The same split preservation
carries a direct-sum null space $\Nsp=(\Nsp\cap\Ssp)\oplus(\Nsp\cap\Tsp)$ to a direct-sum
null space, so the binary certificate $\Pst=\bm0$ is preserved. But $\bm D^{-1}$ does not
preserve inner products, so the magnitude $\lVert\Pst\rVert_{\mathrm F}$ itself may change.
\end{proof}

\noindent Lemma~\ref{lem:precond} lets us report the four integer counts and the
block-diagonality certificate $\Pst=\bm0$ computed on the preconditioned Fisher, secure that
they equal those of the physical, unscaled information. The reported norm
$\lVert\Pst\rVert_{\mathrm F}$ is instead a basis/preconditioner-dependent magnitude
(invariant only under orthogonal basis change, Proposition~\ref{prop:invariance}); for the
anchored networks it is exactly $0$, so no reported number is affected.

\section{The coupled datum-defect classification}
\label{sec:theorem}

This section states and proves the paper's proven contribution: a \emph{basis-invariant}
classification of the null space of the coupled Fisher information \eqref{eq:fisher} into
its spatial, temporal, and genuinely coupled parts; the exact datum gauge of a fully
internal \emph{free-network} (un-anchored, mutual-clock) lunar problem under the SI
convention; the precise criterion that certifies a coupled defect; and the conditions under
which external ties lift each temporal defect---connecting the free-network gauge to the
tied, defect-$0$ realization of Section~\ref{sec:anchor}. Everything here is \Proven{}
algebra; every physical interpretation is labelled where it appears.

\subsection{The coupled datum defect and its projector}
\label{sec:defectdef}

\begin{definition}[Coupled datum defect]
\label{def:defect}
Let $\Nsp=\ker\Imat(\thetav)$ be the null space of the coupled Fisher information, of
dimension $d=\dim\Nsp$ (the \emph{datum defect}), and let $\Umat\in\R^{9\times d}$ have
orthonormal columns spanning $\Nsp$. With the spatial and temporal subspaces
$\Ssp,\Tsp$ of \eqref{eq:subspaces}, write $\Umat_{\Ssp}=\Umat[0{:}7,:]\in\R^{7\times d}$
and $\Umat_{\Tsp}=\Umat[7{:}9,:]\in\R^{2\times d}$ for the row-blocks of $\Umat$ on the two
subspaces. Define
\begin{align}
  \dim(\Nsp\cap\Ssp) &\;=\; d-\rank(\Umat_{\Tsp}) &&\text{(purely spatial defects)},
  \label{eq:dimspatial}\\
  \dim(\Nsp\cap\Tsp) &\;=\; d-\rank(\Umat_{\Ssp}) &&\text{(purely temporal defects)},
  \label{eq:dimtemporal}\\
  \operatorname{coupled\,dim} &\;=\; \rank(\Umat_{\Ssp})+\rank(\Umat_{\Tsp})-d
     &&\text{(coupled defects)}.
  \label{eq:coupleddim}
\end{align}
Equivalently, and manifestly basis-free, let $\Pnull=\Umat\Umat^\top$ be the orthogonal
projector onto $\Nsp$ and let $\Pst=\Pnull[0{:}7,\,7{:}9]\in\R^{7\times2}$ be its
spatial--temporal off-diagonal block. A coupled defect is \emph{certified} by
$\Pst\neq\bm0$; the magnitude $\lVert\Pst\rVert_{\mathrm F}$ is an orthogonally-invariant
(but, by Lemma~\ref{lem:precond}, not preconditioner-invariant) measure of its strength.
\end{definition}

The three counts in \eqref{eq:dimspatial}--\eqref{eq:coupleddim} partition the defect: a
purely spatial null direction is invisible in the temporal rows, a purely temporal one is
invisible in the spatial rows, and a coupled direction is visible in both. Their sum is
$\dim(\Nsp\cap\Ssp)+\dim(\Nsp\cap\Tsp)+\operatorname{coupled\,dim}=d$ \emph{identically},
by the definition of $\operatorname{coupled\,dim}$ in \eqref{eq:coupleddim}; no
``residual'' case exists. That count is moreover nonnegative: the vertical stack
$[\Umat_{\Ssp};\Umat_{\Tsp}]$ of the two row-blocks is $\Umat$ itself, so
$\rank[\Umat_{\Ssp};\Umat_{\Tsp}]=\rank\Umat=d\le\rank(\Umat_{\Ssp})+\rank(\Umat_{\Tsp})$
by rank subadditivity under row stacking, giving
$\operatorname{coupled\,dim}=\rank(\Umat_{\Ssp})+\rank(\Umat_{\Tsp})-d\ge0$. The projector
form is the object the engine reports, and it is the right object because of the following
invariance, which is the paper's central rigor point.

\subsection{Basis invariance}
\label{sec:invariance}

A null space does not come with a canonical basis: any $\Umat\bm Q$ with $\bm Q$ orthogonal
spans the same $\Nsp$, and an eigensolver returns one arbitrary such $\Umat$. A
classification that depended on the individual columns of $\Umat$ would therefore be
ill-posed. The naive quantity ``the energy of a null \emph{vector} in the temporal block,''
$\lVert\bm u_{\Tsp}\rVert$ for a chosen null vector $\bm u$, is exactly such an ill-posed
quantity---it changes under $\Umat\mapsto\Umat\bm Q$. The classification of
Definition~\ref{def:defect} avoids this trap.

\begin{proposition}[Basis invariance of the classification]
\label{prop:invariance}
The three dimensions \eqref{eq:dimspatial}--\eqref{eq:coupleddim} and the coupling strength
$\lVert\Pst\rVert_{\mathrm F}$ depend only on the subspace $\Nsp$, not on the choice of
orthonormal basis $\Umat$. Formally, they are invariant under $\Umat\mapsto\Umat\bm Q$ for
every orthogonal $\bm Q\in\R^{d\times d}$.
\end{proposition}

\begin{proof}
Let $\bm P_{\Ssp}$ and $\bm P_{\Tsp}$ be the fixed coordinate projectors onto $\Ssp$ and
$\Tsp$ (selecting rows $0$--$6$ and $7$--$8$). For the row-block $\Umat_{\Tsp}=\bm P_{\Tsp}\Umat$,
$\rank(\Umat_{\Tsp})=\dim(\bm P_{\Tsp}\Nsp)$, which is a property of the subspace $\Nsp$ and the
fixed projector, independent of the basis; the same holds for $\rank(\Umat_{\Ssp})$. Hence
\eqref{eq:dimspatial}--\eqref{eq:coupleddim}, being integer combinations of these ranks and
$d=\dim\Nsp$, are subspace invariants. For the projector, orthogonal $\bm Q$ gives
$(\Umat\bm Q)(\Umat\bm Q)^\top=\Umat\bm Q\bm Q^\top\Umat^\top=\Umat\Umat^\top=\Pnull$, so
$\Pnull$ and therefore its fixed sub-block $\Pst$ and $\lVert\Pst\rVert_{\mathrm F}$ are
literally unchanged. (Indeed $\Pnull$ is the unique orthogonal projector with range $\Nsp$,
so it is intrinsic to the subspace.)
\end{proof}

\begin{remark}[The projector zeroness certifies coupling; its norm measures strength]
\label{rem:projnorm}
$\lVert\Pst\rVert_{\mathrm F}=0$ if and only if $\Pnull$ is block-diagonal in
$\Ssp\oplus\Tsp$, i.e.\ $\Nsp=(\Nsp\cap\Ssp)\oplus(\Nsp\cap\Tsp)$ is a direct sum of a
spatial and a temporal part with no direction mixing the two. This \emph{binary} certificate
$\Pst=\bm0$ is the invariant object: it is unchanged under both orthogonal basis rotation
(Proposition~\ref{prop:invariance}) and the positive-diagonal preconditioner
(Lemma~\ref{lem:precond}). A single null vector $\bm u=(\bm u_{\Ssp};\bm u_{\Tsp})$ with both
blocks nonzero forces $\Pst\neq\bm0$ and contributes to $\operatorname{coupled\,dim}$. The
\emph{value} $\lVert\Pst\rVert_{\mathrm F}$ is a strength magnitude, invariant under
orthogonal basis change but not under the preconditioner; the engine reports
$\lVert\Pst\rVert_{\mathrm F}=0.5$ for the canonical single coupled direction
$\bm u=(\bm e_{\text{scale}}+\bm e_{\dtau})/\sqrt2$ on an \emph{un-preconditioned} synthetic
matrix, and a unit test confirms this value is unchanged when the null basis is rotated by an
arbitrary angle (Section~\ref{sec:anchor}). For the anchored networks $\Pst=\bm0$, so the
strength value enters no validated number.
\end{remark}

\subsection{The classification theorem and the exact defect count}
\label{sec:classify}

\begin{theorem}[Coupled free-network datum gauge: classification and count]
\label{thm:classify}
Consider a fully internal \emph{free-network} lunar PNT problem: mutual one-way
pseudoranges among orbiters and surface beacons with common-mode clocks and \emph{no}
external position or time tie (model~(A) of Section~\ref{sec:models}), under the SI
convention of Section~\ref{sec:si}, with coupled Fisher \eqref{eq:fisher} assembled from the
free/differential rows in which a common timescale offset and rate cancel \eqref{eq:free}.
Then:
\begin{enumerate}
\item[(i)] \textbf{Exact defect.} The datum gauge defect is
$d=6+2=8$: three translations, three rotations, the clock offset $\dtau$, and the clock
rate $\dalpha$. Scale $s$ is \emph{observable} (not a defect), because under SI a
pseudorange is scale-bearing through the fixed $c$.
\item[(ii)] \textbf{Temporal defects are pure.} In the free/differential design the clock
term enters each one-way range as a difference (receiver minus transmitter), so a common
offset $\dtau$ and a common rate $\dalpha$ cancel in every internal observation and their
datum columns are \emph{identically zero} \eqref{eq:free}; hence $\bm e_{\dtau}$ and
$\bm e_{\dalpha}$ lie in $\ker\Imat$ exactly, giving $\dim(\Nsp\cap\Tsp)=2$ and
$\operatorname{coupled\,dim}=0$ for the free network alone. (In the \emph{anchored} model~(B),
by contrast, the offset column is $+1$---an absolute offset against a known reference
\eqref{eq:oneway}---and is therefore observable; that model is treated in
Section~\ref{sec:anchor}.)
\item[(iii)] \textbf{The joint gauge is not a direct sum in general.} Once an external tie
that couples the subspaces is admitted---in particular the relativistic rate tie
\eqref{eq:ratetie}, whose scale column is nonzero---the null space acquires a direction with
$\Pst\neq\bm0$, so the joint frame$\oplus$timescale gauge is \emph{not} the direct sum of
the spatial gauge of \citep{baweja2026datum} and the temporal gauge.
\end{enumerate}
\end{theorem}

\begin{proof}
\emph{(i), (ii).} Under SI (Section~\ref{sec:si}) the scale column of every one-way row is
the physical range's dependence on the lunar radius through the fixed $c$, which is nonzero
and geometry-dependent across a network of distinct line-of-sight directions; hence scale
is not annihilated and contributes no null direction (this is the P1 result that a
geometrically diverse network observes scale \citep{baweja2026datum}, imported unchanged).
The rigid-body translation ($3$) and rotation ($3$) directions are the classical Helmert
datum defect of a free network: a common translation and a common infinitesimal rotation of
the point set are unobservable to internal ranges \citep{baarda1973,teunissen2000}, giving
six spatial null directions. For the temporal columns, the free/differential design
\eqref{eq:free} carries the clock term as a receiver-minus-transmitter difference, so a
common offset $\dtau$ and a common rate $\dalpha$ leave every internal observation
invariant---their datum columns are \emph{identically zero}---and $\bm e_7=\bm e_{\dtau}$,
$\bm e_8=\bm e_{\dalpha}$ are annihilated by $\Imat$. (This is the defining property of the
free network; it is \emph{not} the anchored row \eqref{eq:oneway}, whose $+1$ offset column
makes the offset observable, since there $(\Imat\bm e_7)_7=\sum_i g_{7,i}^2=N>0$.) Thus
$\ker\Imat\supseteq\operatorname{span}\{\text{6 rigid},\bm e_7,\bm e_8\}$, $d\ge8$; and
because a geometrically diverse free network observes everything else (scale and all
relative geometry), $d=8$ with the temporal pair lying entirely in $\Tsp$, so
$\dim(\Nsp\cap\Tsp)=2$ and $\operatorname{coupled\,dim}=0$.

\emph{(iii).} Add the rate tie \eqref{eq:ratetie}. Its row
$\bm g^{\mathrm{rt}}=-\tfrac{\partial\alpha}{\partial s}\,\bm e_{s}+\bm e_{\dalpha}$ has a
nonzero scale component, so it constrains the combination
$\dalpha-\tfrac{\partial\alpha}{\partial s}\,\delta s$. If scale is otherwise fixed, this
removes the $\dalpha$ defect (below); in the free-network sub-case where no independent
scale tie is present, the surviving null direction is
$\bm u\propto\bigl(\tfrac{\partial\alpha}{\partial s}\,\bm e_{s}+\bm e_{\dalpha}\bigr)$,
which has a nonzero spatial block ($\bm e_s\in\Ssp$) and a nonzero temporal block
($\bm e_{\dalpha}\in\Tsp$); by Remark~\ref{rem:projnorm}, $\Pst\neq\bm0$. Hence the joint
gauge contains a direction mixing $\Ssp$ and $\Tsp$ and is not $\Ssp$-$\Tsp$ block-diagonal.
\end{proof}

\begin{remark}[The off-diagonal criterion: necessary, not sufficient]
\label{rem:necessary}
It is tempting to certify a coupled defect from the \emph{Fisher} rather than its null
space---by checking whether the Fisher's spatial--temporal block $\Imat_{\Ssp\Tsp}$ is
nonzero. This is wrong. $\Imat_{\Ssp\Tsp}\neq\bm0$ (the Fisher is not block-diagonal) is
\emph{necessary} for a coupled defect but \emph{not sufficient}: the coupling can live
entirely among the \emph{observable} directions, leaving the null space a clean direct sum.
Concretely, a network whose scale$\leftrightarrow$offset block is nonzero (they are
correlated) but full-rank has $\Imat_{\Ssp\Tsp}\neq\bm0$ and yet $\operatorname{coupled\,dim}=0$,
$\Pst=\bm0$---the very case of the well-posed anchor network of
Section~\ref{sec:anchor}, which has $\operatorname{coupled\,dim}=0$ despite a
non-block-diagonal Fisher. A coupled \emph{defect} is certified \emph{only} through the null
space: by $\Pst\neq\bm0$, equivalently $\operatorname{coupled\,dim}>0$ via the rank formula
\eqref{eq:coupleddim}.
\end{remark}

\subsection{Lifting the temporal defects}
\label{sec:lift}

The temporal defects are lifted by external ties, and the mechanism is the intersection law
for the null spaces of summed positive-semidefinite (PSD) information blocks.

\begin{lemma}[Intersection law]
\label{lem:intersect}
For PSD matrices $\Imat_A,\Imat_B\succeq0$,
$\ker(\Imat_A+\Imat_B)=\ker\Imat_A\cap\ker\Imat_B$.
\end{lemma}

\begin{proof}
For PSD $\Imat_A,\Imat_B$ and any $\bm v$, $\bm v^\top(\Imat_A+\Imat_B)\bm v=
\bm v^\top\Imat_A\bm v+\bm v^\top\Imat_B\bm v$ is a sum of two nonnegative terms, so it
vanishes iff both vanish, i.e.\ iff $\Imat_A\bm v=\bm0$ and $\Imat_B\bm v=\bm0$ (a PSD matrix
annihilates exactly the vectors of zero quadratic form). Hence
$\ker(\Imat_A+\Imat_B)=\ker\Imat_A\cap\ker\Imat_B$.
\end{proof}

\begin{theorem}[Lift of the free-network temporal defects]
\label{thm:lift}
Let $\Imat$ be the \emph{free-network} Fisher of Theorem~\ref{thm:classify}, whose temporal
gauge is $\ker\Imat\supseteq\operatorname{span}\{\bm e_{\dtau},\bm e_{\dalpha}\}$ (the two
free-network temporal defects). External ties reduce this gauge, connecting the free-network
defect of eight to the fully-tied, defect-$0$ realization of Section~\ref{sec:anchor}:
\begin{enumerate}
\item[(a)] \textbf{Offset.} Adding a time-tie block \eqref{eq:timetie} (rank-one at
$\bm e_{\dtau}$) removes $\bm e_{\dtau}$ from the null space: by Lemma~\ref{lem:intersect} the
combined null space is $\ker\Imat\cap\ker\Imat_{\mathrm{tt}}$, and $\Imat_{\mathrm{tt}}$
annihilates every vector except those with a $\dtau$ component, so $\bm e_{\dtau}$ is lifted
while $\bm e_{\dalpha}$ survives.
\item[(b)] \textbf{Rate.} Adding the rate tie \eqref{eq:ratetie} constrains
$\dalpha-\tfrac{\partial\alpha}{\partial s}\,\delta s$. Consider the state \emph{after} the
offset has been tied (part~a) and with scale observable (as under SI) or independently tied:
the surviving null space then has no scale component ($\delta s=0$) and its only temporal
direction is $\bm e_{\dalpha}$, so the rate tie forces $\dalpha=0$ and lifts it fully. Were
scale instead left free---a free-network sub-case, not the SI default---the rate tie alone
would fuse the rate and scale directions into a single coupled scale$\leftrightarrow$rate
null direction (reducing that sector's dimension by one) rather than remove the defect. Thus
(a) followed by (b) lifts \emph{both} free-network temporal defects, and the fully-tied
network reaches defect $0$ (Section~\ref{sec:anchor}).
\end{enumerate}
\end{theorem}

\begin{proof}
Both parts are Lemma~\ref{lem:intersect} applied to $\Imat+\Imat_{\mathrm{tie}}$. For (a),
$\Imat_{\mathrm{tt}}=\bm e_{\dtau}\bm e_{\dtau}^\top$ has kernel
$\{\bm v:v_{\dtau}=0\}$; intersecting with $\ker\Imat$ removes exactly $\bm e_{\dtau}$. For
(b), $\Imat_{\mathrm{rt}}=\bm g^{\mathrm{rt}}(\bm g^{\mathrm{rt}})^\top$ has kernel
$\{\bm v:\bm g^{\mathrm{rt}}\!\cdot\!\bm v=0\}=\{\bm v:v_{\dalpha}=\tfrac{\partial\alpha}
{\partial s}v_{s}\}$; intersecting with $\ker\Imat$, and using that under SI $\ker\Imat$ has
no scale component ($v_s=0$ on the null space), forces $v_{\dalpha}=0$, removing
$\bm e_{\dalpha}$. Without an independent scale tie, $v_s$ is free and the surviving
direction is $\tfrac{\partial\alpha}{\partial s}\bm e_s+\bm e_{\dalpha}$, a single coupled
direction.
\end{proof}

\begin{remark}[Exact-null lift versus near-defect lift]
\label{rem:distinct}
The lifts of Theorem~\ref{thm:lift} act on \emph{exact} null directions (defects). They are
distinct from, and must not be conflated with, the raising of a small-but-positive
eigenvalue of the $\{s,\dtau\}$ marginal by additional observations---the
\emph{near}-defect of coupling~1 (Section~\ref{sec:coupling1}), which is never a member of
$\ker\Imat$ in a diverse network. One is a change of the defect $d$; the other is a change
of a strictly positive eigenvalue.
\end{remark}

\section{The two frame--time couplings}
\label{sec:couplings}

Theorem~\ref{thm:classify} tells us \emph{which} directions are undetermined; this section
quantifies \emph{how strongly} the frame and the timescale couple along two physically
distinct channels. The classical roots of both channels are \emph{cited}, not claimed; the
theorem structure around them is \Proven{}; and the magnitudes are \Modelled{}
(first-principles, cross-checked, not certified) and tagged \modtag{}. One of the two
subsections corrects a claim whose naive form is false, and we present the correction
openly as the honest result.

\subsection{Coupling 1: range \texorpdfstring{$\leftrightarrow$}{<->} clock offset}
\label{sec:coupling1}

\paragraph{The unit tie and the classical root.}
For one-way ranging, a radial position shift $\delta d$ and a clock offset $c\,\dtau$ enter
every measurement with the same sign: a range is lengthened equally by moving the target
away or by advancing the clock. The exact exchange rate is $c$,
\begin{equation}
  1\,\text{m}\;\longleftrightarrow\;\frac{1}{c}\;=\;3.336\,\text{ns},
  \label{eq:unittie}
\end{equation}
computed from the engine's $c$ (\code{lunar\_time::C2\_M2\_S2}). This is precisely the
GNSS radial-position/receiver-clock degeneracy, classical in the integrity and positioning
literature \citep{sturza1988,brown1992,parkinson1988}; we cite it and claim no novelty. Its
importance here is that it is the strongest frame--time coupling, and the natural place to
look for a \emph{coupled datum defect}.

\paragraph{It is a near-defect, not an exact one.}
Crucially, \eqref{eq:unittie} is a \emph{near}-defect: for a single user with poor local
vertical geometry the scale/radial and offset signatures are nearly parallel, giving a
small---but strictly positive---eigenvalue, not a member of $\ker\Imat$. It becomes an
exact null direction only in the degenerate limit of perfectly co-directional geometry.
Accordingly it does \emph{not} contribute to the defect $d$ of Theorem~\ref{thm:classify};
we quantify it instead by the $2\times2$ Schur-complement (marginal) Fisher onto
$\{s=\text{idx }3,\ \dtau=\text{idx }7\}$, eliminating the other seven parameters
(\code{coupled\_marginal\_fisher} / \code{coupled\_marginal\_eigs}; the algebra is in
Appendix~\ref{app:methods}). The smallest eigenvalue $\lambda_{\min}$ of this marginal is
the coupling-1 near-defect strength: the degree to which the lunar body's scale
(radial breathing) is observable \emph{independently} of the common clock offset in a given
one-way network.

\paragraph{The corrected, defensible finding.}
The tempting design claim is: ``two-way ranging (which cancels the node clock, cf.\
\eqref{eq:twoway}) selectively breaks the scale$\leftrightarrow$offset trade and lifts the
near-defect, so a lunar system needs two-way ranging to separate frame scale from
timescale offset.'' On a real-code-path probe this claim proves \emph{false at the datum
level}, and we report the correction. On a geometrically diverse multi-beacon
multi-epoch one-way network (five beacons $\times$ six orbiters $\times$ four epochs,
$120$ one-way rows, real DE440 orientation), the coupled Fisher is \emph{full-rank}
($d=0$) and the $\{s,\dtau\}$ marginal is positive-definite: its eigenvalues are
$[\,1.872,\ 36.00\,]$ (ascending), both strictly positive. Scale is already observable
independently of the offset. We then compare, at equal observation count, two ways of
adding $120$ rows to this base (Table~\ref{tab:coupling1}): (a) $120$ two-way rows at the
same geometry, and (b) $120$ \emph{more} one-way rows at a jittered geometry (argument of
latitude shifted by $7^\circ$). Both raise $\lambda_{\min}$ of the marginal; the ratio of
the two lifts is
\begin{equation}
  \frac{\text{lift}_{\text{two-way}}}{\text{lift}_{\text{more one-way}}}\;\approx\;1.00,
  \label{eq:liftratio}
\end{equation}
altitude-independent within the tested networks. In words: \textbf{in a geometrically
diverse network, two-way ranging provides no selective scale$\leftrightarrow$offset
decoupling beyond data volume.} The strong $1\,\text{m}\leftrightarrow3.336\,\text{ns}$
degeneracy is a \emph{per-node}, single-user, poor-local-geometry effect (the classical
GNSS radial/clock degeneracy) that does \emph{not} aggregate into a datum-level near-defect
once the beacon and orbiter directions are geometrically diverse.

\begin{table}[t]
\centering
\caption{Coupling~1 at the datum level (\Modelled{} geometry, real DE440 orientation;
\modtag{}). On a geometrically diverse $120$-row one-way network the coupled Fisher is
full-rank and the $\{s,\dtau\}$ Schur marginal is positive-definite. Adding $120$ two-way
rows raises the marginal $\lambda_{\min}$ by the \emph{same} factor as adding $120$ more
one-way rows: two-way ranging gives no selective decoupling beyond data volume. The
marginal eigenvalue pair $[1.872,\,36.00]$ is reproduced by the independent numpy/LAPACK
oracle (Section~\ref{sec:anchor}).}
\label{tab:coupling1}
\small
\begin{tabular}{@{}lcc@{}}
\toprule
\textbf{Network} & $\{s,\dtau\}$ marginal $\lambda_{\min}$ & \textbf{status} \\
\midrule
Base: $120$ one-way rows (diverse, defect $0$) & $1.872$ & \Validated{} (Table~\ref{tab:anchor}) \\
$+\,120$ two-way rows (same geometry)          & lift $>1$ & \Modelled{} \\
$+\,120$ more one-way rows (jittered $7^\circ$) & lift $>1$ & \Modelled{} \\
lift ratio (two-way / more one-way)             & $\approx1.00$ & \Modelled{} \emph{(the finding)} \\
\bottomrule
\end{tabular}
\end{table}

\begin{remark}[The idealized single-node regime, for contrast]
\label{rem:singlenode}
The naive claim is not \emph{empty}: in an idealized single-node, near-rank-one regime
(one beacon, poor vertical geometry) a two-way tie \emph{does} raise $\lambda_{\min}$ by
orders of magnitude, because there the scale and offset signatures really are nearly
parallel. That regime is a unit-test fixture, not a lunar constellation. The design-relevant
statement is the diverse-network one \eqref{eq:liftratio}; we keep the two regimes strictly
separate, and neither is a coupled \emph{defect} (Remark~\ref{rem:distinct}).
\end{remark}

This yields the first design law, developed in Section~\ref{sec:designlaws}:
\emph{geometric diversity separates the lunar frame scale from the timescale offset in
one-way ranging; a two-way tie or an external time tie is required only in the
low-diversity / single-user regime.}

\subsection{Coupling 2: relativistic rate \texorpdfstring{$\leftrightarrow$}{<->} frame}
\label{sec:coupling2}

The second channel is the dependence of a lunar-surface clock's \emph{rate} on the
frame it is realized in. A clock runs at a rate set by the gravitational potential and the
squared speed of its location (relativistic geodesy and the IAU relativistic timescale
framework \citep{soffel2003iau,mueller2018clocks}); perturbing the frame datum---its scale,
a radial position, or an inferred velocity---perturbs that rate. This is a real coupling,
and we ask only how \emph{strong} it is.

\paragraph{The frame\,$\to$\,rate Jacobian (\Modelled{}, \modtag{}).}
Writing $\Umoon=GM_{\mathrm{M}}/\rmoon$ for the lunar surface potential and
$\gmoon=GM_{\mathrm{M}}/\rmoon^2$ for the surface gravity, the engine's
\code{rate\_frame\_jacobian} returns the three first-principles coefficients
\begin{align}
  \frac{\partial\alpha}{\partial s}
    &\;=\; +\frac{\Umoon}{c^2}
    \;=\; +3.140\times10^{-11}\ \text{per unit fractional scale},
    \label{eq:dads}\\
  \frac{\partial\alpha}{\partial v}
    &\;=\; -\frac{|v|}{c^2}
    \;\approx\; -1.11\times10^{-14}\ \text{per (m/s)},
    \label{eq:dadv}\\
  \frac{\partial\alpha}{\partial r_{\mathrm{radial}}}
    &\;=\; +\frac{\gmoon}{c^2}
    \;=\; +1.807\times10^{-17}\ \text{per metre}.
    \label{eq:dadr}
\end{align}
The radial and scale coefficients satisfy the exact identity
$\partial\alpha/\partial r=(\partial\alpha/\partial s)/\rmoon$ (unit-tested to
$10^{-12}$), since $\gmoon/c^2=(\Umoon/\rmoon)/c^2$. The velocity coefficient is
epoch-dependent because the geocentric Moon speed varies by several percent over the lunar
month (the $969$ and $997$~m/s values below already differ by nearly $3\%$); the engine
prints $-1.078\times10^{-14}$ at the fixture epoch 2024-01-01
($|v|\approx969$~m/s) versus $\approx-1.11\times10^{-14}$ near J2000
($|v|\approx997$~m/s). We quote both as honest epoch variation, not an error. The signs
combine as
\begin{equation}
  \dalpha
  \;=\; \frac{\partial\alpha}{\partial s}\,\delta s
  \;+\; \frac{\partial\alpha}{\partial r}\,\delta r
  \;-\; \Bigl|\frac{\partial\alpha}{\partial v}\Bigr|\,|\delta v|,
  \label{eq:dalphacombine}
\end{equation}
the velocity term subtracting (a faster clock runs slower). The velocity and radial
couplings are physical but, as noted in Section~\ref{sec:rows}, do not map onto the
nine-vector datum; only the scale coupling appears as the datum-level rate tie
\eqref{eq:ratetie}.

\paragraph{Propagating an ILRF-class frame datum to the rate (\Modelled{}, \modtag{}).}
Take frame-tie uncertainties at the level a lunar analogue of the terrestrial reference
frame realization would achieve: a scale tie $\delta s\approx6.9\times10^{-8}$
(i.e.\ $0.12$~m of radial equivalent, $0.12/\rmoon$), a velocity of
$\delta v\lesssim1$~mm/s, and a radial position of $\delta r\lesssim0.1$~m. Through
\eqref{eq:dads}--\eqref{eq:dadr} these propagate to
\begin{equation}
  \dalpha(\text{scale})\approx2.17\times10^{-18},\quad
  \dalpha(v,\,1\,\text{mm/s})\approx1.11\times10^{-17},\quad
  \dalpha(\text{radial},\,0.1\,\text{m})\approx1.81\times10^{-18},
  \label{eq:propagate}
\end{equation}
so the frame \emph{realization} perturbs the lunar timescale rate by at most a
frame-realization envelope
\begin{equation}
  \boxed{\;\dalpha_{\text{frame realization}}\;\lesssim\;\sim2\times10^{-17}\;}
  \qquad\text{(realization envelope, kinetic-dominated).}
  \label{eq:framebound}
\end{equation}
These two budgets must be kept apart. Only the scale term of \eqref{eq:propagate} is a
genuine \emph{datum}-channel gain: $\dalpha(\text{scale})\approx2.2\times10^{-18}$ is what a
nine-vector datum-defect $\delta s$ propagates to the rate (Section~\ref{sec:designlaws}).
The larger velocity and radial contributions are frame-\emph{realization} /
orbit-determination errors that, as \eqref{eq:ratetie} and Section~\ref{sec:rows} note, do
\emph{not} map onto the datum; the envelope \eqref{eq:framebound} merely absorbs them for a
conservative bound. The gauge-robustness conclusion holds \emph{a fortiori}---the true
datum-channel gain is an order of magnitude below the realization envelope. Nothing larger
exists in either budget: the tidal-gradient contribution is $\sim2.7\times10^{-22}$ per
metre, and the external/tidal potential is excluded from $\Umoon$ by the equivalence
principle (a freely-falling body does not sense the uniform external potential, so we do not
add $GM_{\oplus}/d$). We quote the envelope as $\lesssim\!\sim\!2\times10^{-17}$, not
$\lesssim10^{-17}$, to keep the kinetic term honest.

\paragraph{The reference-surface definition band (\Modelled{}, \modtag{}).}
Against this estimation coupling stands the \emph{definitional} freedom of the lunar
timescale rate---the choice of reference surfaces, reference potential $W_0$, and averaging
window that a standard must fix. The lunar-clock-rate band relative to TT spans
$[56,59]$~$\mu$s/day, consistent with the lunar-timescale figures discussed in the Moonlight
reference-system literature
% confirm the [56,59] $\mu$s/day figure against fienga2024moonlight at camera-ready
\citep{fienga2024moonlight}; its half-width $\pm1.5$~$\mu$s/day
converts to a dimensionless rate
\begin{equation}
  \pm\frac{1.5\times10^{-6}\ \text{s/day}}{86400\ \text{s/day}}
  \;=\;\pm1.736\times10^{-11},
  \label{eq:band}
\end{equation}
about $10^{6}$ times the estimation coupling \eqref{eq:framebound}. The two axes are
orthogonal: one is set by \emph{how well the frame is estimated}, the other by \emph{how the
reference surfaces are defined}. The band's physical content includes named contributors
inside the lunar potential---tidal deformation $\approx2.4\times10^{-16}$ and $J_2$
oblateness $\approx6\times10^{-15}$, together with the periodic solar tidal contribution to
the Earth--Moon system $\approx6\times10^{-14}$ (periodic, and \emph{larger} than the $J_2$
term)---this named list is illustrative, not exhaustive (dimensionless; physics-literature
reference values, not engine-computed); its full $\pm1.7\times10^{-11}$ width is dominated by
the definitional $W_0$/selenoid-surface choice and the averaging window. A model-center
cross-check reconciles the band midpoint from first principles: the Earth-geoid redshift
that defines TT contributes $\approx60.2$~$\mu$s/day, the lunar surface potential subtracts
$\approx2.71$~$\mu$s/day, and the lunar orbital kinetic term subtracts
$\approx0.48$~$\mu$s/day, giving $60.2-2.71-0.48=57.0$~$\mu$s/day, comfortably within
$[56,59]$ (all \Modelled{}, first-principles).

\paragraph{Where the estimation coupling would matter.}
The estimation bound \eqref{eq:framebound} of $\sim2\times10^{-17}$ sits below the
stability of near-term deployed optical clocks ($\sim10^{-16}$) and would become relevant
only for laboratory-grade clocks at $\sim10^{-18}$. It is, for any foreseeable operational
lunar timescale, negligible against the definitional band.

\paragraph{Headline (\Proven{} logic, \Modelled{} magnitudes).}
The lunar timescale rate is \textbf{gauge-robust to the frame estimation datum}
($\lesssim\!\sim\!2\times10^{-17}$) but \textbf{gauge-sensitive to the reference-surface
definition} ($\pm1.7\times10^{-11}$). The dominant gauge of the rate is
\emph{definitional}, not \emph{estimational}: no PNT observation reduces it, and a standard
must \emph{fix} the surfaces and conventions. This is the time-analogue of the spatial
result of \citep{baweja2026datum}---the datum is fixed by convention and external ties, not
estimated---transported to the rate of the timescale.

\section{The real-data validated anchor}
\label{sec:anchor}

The paper's single externally \Validated{} claim is narrow and precise: that the
coupled-Fisher \emph{linear-algebra pipeline}---preconditioned assembly
\eqref{eq:fisher}, eigendecomposition, the basis-invariant null-space classification of
Definition~\ref{def:defect}, and the $\{s,\dtau\}$ Schur marginal of
Section~\ref{sec:coupling1}---reproduces an \emph{independent} numpy/LAPACK computation of
the same pipeline, to relative and absolute error below $10^{-3}$, on observation rows
whose spatial partials carry the \emph{real} DE440 lunar physical-axis orientation
\citep{park2021de440}. This section states what is validated and, just as carefully, what
is not.

\subsection{Two anchored networks}
\label{sec:networks}

Both networks are built by the generator \code{gen\_coupled\_gauge\_rows}, which forms
each one-way row \eqref{eq:oneway} through the real DE440 lunar physical-axis orientation
path: the row routine \code{orbiter\_range\_row\_datum7} internally calls
\code{de440\_moon\_pa\_body\_to\_inertial}, so columns $0$--$6$ of every row carry the real
physical libration at the epoch, not an approximation. The geometry (five near-side beacons,
six circular orbiters at $2000$~km, up to four two-day epochs from 2024-01-01 TT) is
\Modelled{}; the DE440 orientation is the external ingredient.

\begin{table}[t]
\centering
\caption{The \Validated{} anchor: the coupled-Fisher pipeline on real DE440-orientation
rows, reproduced by an independent numpy/LAPACK oracle to relative and absolute error below
$10^{-3}$ (integer counts exact). The \code{well\_posed} network is full-rank (defect $0$)
with a positive-definite $\{s,\dtau\}$ marginal $[1.872,\,36.00]$; the \code{single\_epoch}
network ($t_{\mathrm e}=0$, rate column identically zero) has an \emph{exact} rate-direction
defect (defect $1$, $\dim(\Nsp\cap\Tsp)=1$). Both have $\operatorname{coupled\,dim}=0$ and
$\lVert\Pst\rVert_{\mathrm F}=0$: the anchored branch does \emph{not} exercise nonzero
coupling.}
\label{tab:anchor}
\footnotesize
\setlength{\tabcolsep}{5pt}
\begin{tabular}{@{}lcccccc@{}}
\toprule
\textbf{Network} & \textbf{rows} & $d$ & $\dim(\Nsp\!\cap\!\Ssp)$ & $\dim(\Nsp\!\cap\!\Tsp)$ &
$\operatorname{cpld.}$ & $\{s,\dtau\}$ marg.\ eigs \\
\midrule
\code{well\_posed}   & $120$ & $0$ & $0$ & $0$ & $0$ & $[1.8720,\ 35.9995]$ \\
\code{single\_epoch} & $30$  & $1$ & $0$ & $1$ & $0$ & $[0.4171,\ 5.9709]$ \\
\bottomrule
\end{tabular}\\[3pt]
{\footnotesize $d$: defect; $\dim(\Nsp\!\cap\!\Ssp)$/$\dim(\Nsp\!\cap\!\Tsp)$:
purely spatial/temporal defects; $\operatorname{cpld.}$: coupled dimension.}
\end{table}

\paragraph{\code{well\_posed} ($120$ one-way rows).}
Five beacons $\times$ six orbiters $\times$ four epochs. The full $9\times9$ Fisher is
full-rank: all nine eigenvalues are strictly positive,
\begin{equation*}
  [\,0.5685,\ 0.7266,\ 1.0761,\ 1.7837,\ 6.2954,\ 48.228,\ 54.055,\ 79.058,\ 164.95\,],
\end{equation*}
so defect $d=0$ and there is no datum-level near-defect. The $\{s,\dtau\}$ Schur marginal is
positive-definite with eigenvalues $[\,1.8720,\,35.9995\,]$ (ascending): scale is observable
independently of the clock offset. This is the base network of Table~\ref{tab:coupling1} and
the concrete witness for Remark~\ref{rem:necessary}---a non-block-diagonal Fisher with
$\operatorname{coupled\,dim}=0$.

\paragraph{\code{single\_epoch} ($30$ one-way rows).}
The same geometry at one epoch with elapsed time $t_{\mathrm e}=0$, which makes the rate
column $t_{\mathrm e}/T_{\mathrm B}$ identically zero in every row and collapses the rate
column of the Fisher. The defect is exactly $1$, entirely temporal
($\dim(\Nsp\cap\Tsp)=1$, $\operatorname{coupled\,dim}=0$): the rate $\dalpha$ is
unobservable. This is an \emph{exact analytic property} of a zero column, not an empirical
near-defect, and we label it so. The smallest \emph{nonzero} eigenvalue is $0.1605$, more
than six orders of magnitude above the numerical rank threshold
$\text{rel\_tol}\cdot\lambda_{\max}=10^{-9}\times44.97\approx4.5\times10^{-8}$---a clean
spectral gap, so the defect count is unambiguous.

\paragraph{Why the anchor has defect $0$, not $8$.}
Both networks realize model~(B) of Section~\ref{sec:models}: their orbiters sit at known
inertial DE440 ephemeris positions with a known time reference---an implicit \emph{absolute}
frame and time tie---so each one-way row carries the $+1$ offset column of \eqref{eq:oneway}
and the datum is observable. Consequently the \code{well\_posed} network has defect $0$, not
the eight of the fully-internal free-network gauge of Theorem~\ref{thm:classify}: the
anchoring is precisely the external tie that Theorem~\ref{thm:lift} shows lifts the six rigid
directions, the offset, and---given the observable scale and a nonzero time baseline---the
rate. The \code{single\_epoch} network retains a single defect, the rate, not for want of
anchoring but because its zero time baseline ($t_{\mathrm e}=0$) makes the rate column
identically zero, leaving $\dalpha$ unobservable regardless of the tie. The anchor is thus
the fully-tied endpoint of the free-network$\,\to\,$tied arc of Section~\ref{sec:theorem}, so
the theorem and the validated table are consistent, not contradictory.

\subsection{The independent oracle and non-circularity}
\label{sec:oracle}

The reference values are produced by \code{scripts/gen\_coupled\_gauge\_ref.py}, which reads
the \emph{rows} from the vendored \code{network.json}, applies the identical
preconditioning (divide columns $3$--$6$ by $\rmoon=1737400$~m), and computes the Fisher,
its eigenvalues (\code{np.linalg.eigvalsh}), the null-space classification
(\code{np.linalg.eigh} on the Gram sub-blocks), and the Schur marginal
(\code{np.linalg.pinv}) through LAPACK---a different eigensolver (divide-and-conquer) and a
different pseudo-inverse (SVD) from the engine's Jacobi eigensolver and spectral-sum
pseudo-inverse. The comparison is \emph{non-circular}: \code{reference.json} is derived from
the input rows, not from any Rust output. Byte-consistency is enforced by rounding every row
float to six decimals \emph{before} both writing \code{network.json} and reading it in the
oracle, so both implementations consume the same IEEE-754 doubles and the test isolates the
two \emph{algorithms}. The Rust integration test
\code{tests/lunar\_coupled\_gauge\_reference.rs} asserts agreement to relative and absolute
error below $10^{-3}$ (integer counts exact) or fails the build.

\subsection{The honesty boundary: what is and is not validated}
\label{sec:boundary}

We state the scope with the discipline the paper is built on.

\begin{itemize}
  \item \textbf{Validated.} The linear algebra---assembly, eigenvalues, the
  defect/subspace/coupled counts, and the $\{s,\dtau\}$ marginal---on real
  DE440-orientation inputs, for the defect-$0$ (\code{well\_posed}) and the
  exact-rate-defect (\code{single\_epoch}) cases, agreeing with an independent numpy/LAPACK
  oracle to $<10^{-3}$.

  \item \textbf{Modelled, not validated.} The network geometry (beacons, orbits, epochs) is
  a representative construction. The coupling magnitudes of Section~\ref{sec:couplings} are
  first-principles engineering bounds cross-checked to the published band, not certified.

  \item \textbf{Proven, not externally validated.} Both anchored networks have
  $\operatorname{coupled\,dim}=0$ and $\lVert\Pst\rVert_{\mathrm F}=0$, so this external row
  does \emph{not} exercise the nonzero-coupling branch of the classifier. That branch---a
  coupled direction with $\Pst\neq\bm0$---is \Proven{} by construction (the algebra of
  Definition~\ref{def:defect} and Proposition~\ref{prop:invariance}) and is covered by
  module unit tests: a direct-sum case (defect $2$, $\operatorname{coupled\,dim}=0$,
  $\lVert\Pst\rVert_{\mathrm F}=0$), a coupled case (defect $1$,
  $\operatorname{coupled\,dim}=1$, $\lVert\Pst\rVert_{\mathrm F}=0.5$), and a
  \emph{rotated-basis invariance} test that rotates the null basis by an arbitrary angle and
  confirms the classification is unchanged (the concrete witness for
  Proposition~\ref{prop:invariance}). We do \emph{not} claim the coupled branch as
  externally validated; it is proven and unit-tested.
\end{itemize}

This is the same firewall discipline as the companion papers
\citep{baweja2026datum,baweja2026commonmode}: exactly one externally validated row, its
scope stated precisely, everything else cited, proven, or modelled.

\section{Two design laws and the reference-system theorem}
\label{sec:designlaws}

The classification (Section~\ref{sec:theorem}) and the two couplings
(Section~\ref{sec:couplings}) yield two design laws for a lunar reference-system standard
and a standard-level theorem that combines them. The logic is \Proven{}; the magnitudes it
rests on are the \Modelled{} couplings, labelled where they enter.

\begin{quote}
\textbf{Design law~1 (frame scale vs.\ timescale offset).} \emph{Geometric diversity
separates the lunar frame scale from the timescale offset in one-way ranging.} A
geometrically diverse multi-beacon multi-epoch one-way network makes the coupled Fisher
full-rank and the $\{s,\dtau\}$ marginal positive-definite (Section~\ref{sec:coupling1},
Table~\ref{tab:anchor}); a two-way tie or an external time tie is required \emph{only} in
the low-diversity / single-user regime, where the classical
$1\,\text{m}\leftrightarrow3.336\,\text{ns}$ near-defect is real. Two-way ranging confers no
selective datum-level advantage over an equal volume of well-distributed one-way
observations.
\end{quote}

\begin{quote}
\textbf{Design law~2 (the rate's residual gauge is definitional).} \emph{The lunar
timescale rate's residual gauge is definitional, not estimational.} The frame estimation
datum perturbs the rate by at most $\sim2\times10^{-17}$
\eqref{eq:framebound}, six orders of magnitude below the $\pm1.7\times10^{-11}$
reference-surface definitional band \eqref{eq:band}. No PNT observation reduces the
definitional freedom; a standard must \emph{fix} the reference surfaces (the potential
$W_0$, the selenoid) and the averaging window. Because the frame estimation barely perturbs
the rate, the rate is gauge-robust to how well the frame is measured and gauge-sensitive to
how the surfaces are defined.
\end{quote}

\begin{theorem}[A lunar reference \emph{system} must fix both a datum and a convention]
\label{thm:system}
A complete lunar reference-system realization must fix \emph{both}: (a)~a spatial datum---%
the six rigid-body gauge directions of Theorem~\ref{thm:classify}, pinned by external ties
as in \citep{baweja2026datum}; and (b)~a temporal convention---the offset $\dtau$, pinned by
an external time tie (Theorem~\ref{thm:lift}a), and the reference surfaces and averaging
window that fix the rate (Design law~2). Neither suffices alone:
\begin{enumerate}
\item[(i)] geometry alone (one-way ranging) cannot separate radial position from time
offset in the low-diversity regime (Section~\ref{sec:coupling1}, coupling~1);
\item[(ii)] the relativistic model alone cannot bootstrap an absolute rate without a spatial
datum---the rate tie \eqref{eq:ratetie} lifts $\dalpha$ only \emph{given} an independent
scale tie (Theorem~\ref{thm:lift}b); and
\item[(iii)] the rate's residual gauge is definitional, not observable, so no amount of data
removes it (Design law~2).
\end{enumerate}
\end{theorem}

\begin{proof}
Each clause restates a proven result. (i) is coupling~1: in the low-diversity regime the
$\{s,\dtau\}$ marginal has a small eigenvalue, so scale/radial and offset are not separated
by one-way geometry alone. (ii) is Theorem~\ref{thm:lift}b: the rate tie constrains only the
combination $\dalpha-\tfrac{\partial\alpha}{\partial s}\delta s$, which fixes $\dalpha$ iff
$\delta s$ is independently fixed. (iii) is Design law~2 with the bound \eqref{eq:framebound}
against the band \eqref{eq:band}. A realization missing (a) has an unobservable
frame; one missing the offset tie has an unobservable $\dtau$ (Theorem~\ref{thm:classify}ii);
one missing the surface/window convention has an undefined rate. Hence both a spatial datum
and a temporal convention are necessary.
\end{proof}

\paragraph{The P1\,$\to$\,P4 propagation.}
Theorem~\ref{thm:system} makes explicit how the spatial datum of \citep{baweja2026datum}
propagates into the timescale. The spatial datum defect enters the timescale rate through
the rate-tie Jacobian $\partial\alpha/\partial s$ \eqref{eq:dads}, so a residual scale-datum
error $\delta s$ maps to a rate error $(\partial\alpha/\partial s)\,\delta s$. This
\emph{datum-channel} gain---the genuine P1$\to$P4 gauge propagation---is
$\approx2.2\times10^{-18}$ for ILRF-class ties (the scale term of \eqref{eq:propagate}). It
must not be conflated with the broader frame-realization envelope
$\lesssim\!\sim\!2\times10^{-17}$ \eqref{eq:framebound}, which additionally absorbs non-datum
velocity and radial orbit-determination errors that Section~\ref{sec:couplings} shows do not
map onto the nine-vector datum. The spatial datum thus \emph{does} propagate into the
rate---closing the P1$\to$P4 loop---yet the rate is gauge-robust to it, and \emph{a
fortiori}: the true datum-channel gain is an order of magnitude below even that envelope. The
joint gauge is not the direct sum of the spatial and temporal gauges
(Theorem~\ref{thm:classify}iii), but along the one channel that couples them the coupling is
numerically negligible against the definitional band. This is the precise, quantified sense
in which the frame and the timescale are ``coupled but weakly so'': coupled in structure
(a nonzero $\Pst$ once the rate tie is admitted), weak in magnitude (six or more orders below
the definitional gauge).

\section{Discussion}
\label{sec:discussion}

\subsection{Relation to the companion arc}
\label{sec:companions}

This paper is the temporal capstone of a four-part arc built on one open engine. P1
\citep{baweja2026datum} classifies the \emph{spatial} datum defect of a lunar network and
the observation geometry (libration, an absolute tie) that resolves it; its Helmert-seven
datum is the spatial block $\Ssp$ here, imported unchanged, including the SI result that a
geometrically diverse network observes scale. P2 \citep{baweja2026interop} quantifies the
irreducible cross-provider interoperability floor---the inter-ephemeris orientation and
position disagreement between independent authoritative realizations. P3
\citep{baweja2026commonmode} shows that the rigid-body frame freedom is exactly the
error class a snapshot RAIM/ARAIM residual test cannot see (the datum defect is a subset of
the integrity blind subspace). The present paper (P4) adds the missing dimension---%
\emph{time}---and the \emph{coupling}: it extends the datum from the spatial seven to the
coupled nine, classifies the joint null space in basis-invariant form, and shows the joint
gauge is not the direct sum of the spatial and temporal gauges. The four share the same
firewall (exactly one externally validated row per paper) and the same open substrate.

\subsection{Implications for LunaNet, Moonlight/LCNS, and the IAU lunar reference system}
\label{sec:implications}

The two design laws translate directly into specification guidance, stated as \Modelled{}
guidance rather than a certified budget.

\emph{On the frame/offset separation.} A LunaNet or Moonlight/LCNS realization
\citep{israel2020lunanet,fienga2024moonlight,lnis2023} that fields a geometrically diverse
beacon and orbiter network can separate the lunar frame scale from the timescale offset with
one-way ranging alone (Design law~1); it need not mandate two-way ranging \emph{for datum
separation} (two-way ranging remains valuable for other reasons---absolute range, clock
synchronization---outside this analysis). The requirement to add a two-way or external time
tie applies specifically to the low-diversity / single-user regime, where the classical
radial/clock degeneracy bites.

\emph{On the timescale rate.} The IAU lunar-reference-system and lunar-coordinate-time
effort, and the Moonlight lunar-time work \citep{fienga2024moonlight}, must fix the rate by
\emph{convention}: the reference potential $W_0$, the selenoid surface, and the averaging
window (Design law~2). This is not a shortfall of the PNT system---no observation can reduce
the definitional freedom---but a statement about what a standard is \emph{for}. The
consolation is that the frame estimation, however imperfect, perturbs the rate by
$\lesssim\!\sim\!2\times10^{-17}$, negligible against the definitional band and below
near-term deployed optical-clock stability; the rate is robust to estimation error and needs
only a defined convention.

\emph{On the joint realization.} Theorem~\ref{thm:system} is the compact requirement: a
lunar reference \emph{system} standard must fix both a spatial datum and a temporal
convention; fixing one and estimating the other leaves an unobservable or undefined gauge.
This is the coupled-gauge counterpart of the spatial design law of \citep{baweja2026datum}
and the interoperability design law of \citep{baweja2026interop}.

\subsection{Limitations}
\label{sec:limitations}

The firewall of Section~\ref{sec:anchor} is also a statement of limitations; we make them
explicit.

\begin{itemize}
  \item \textbf{Modelled geometry.} The beacon and orbiter geometry is a representative
  construction, not a flown constellation. The classification (Theorem~\ref{thm:classify})
  and its counts are algebraic and hold for any network of the stated form; the specific
  eigenvalues and marginal values (Tables~\ref{tab:coupling1}--\ref{tab:anchor}) depend on
  the modelled geometry and would shift for a real network.

  \item \textbf{First-order linearization.} The datum, the observation rows, and the Fisher
  are first-order about a nominal. Second-order geometric effects (e.g.\ the finite-rotation
  residual, small on the P3 nanoradian scale) and non-Gaussian error structure are outside
  the model.

  \item \textbf{Coupling magnitudes are order-of-magnitude engineering bounds.} The three
  Jacobian entries \eqref{eq:dads}--\eqref{eq:dadr} and the propagated bound
  \eqref{eq:framebound} are first-principles, cross-checked to the published band, and
  unit-tested for internal consistency---but they are \emph{not} certified budgets. The
  definitional-band decomposition (tidal, $J_2$, $W_0$/selenoid, window) uses
  physics-literature reference values for the named sub-terms.

  \item \textbf{The coupled branch is proven, not externally validated.} The single
  external row (Section~\ref{sec:anchor}) exercises defect-$0$ and exact-rate-defect
  networks, both with $\operatorname{coupled\,dim}=0$. The nonzero-coupling classification is
  proven by construction and unit-tested (including rotated-basis invariance), not confirmed
  against an external dataset.

  \item \textbf{No flight data, no certification, no TRL claim.} Nothing here is a certified
  timescale realization, a certified frame, or a statement above analysis. There is no
  single-number operational rate claim; the rate is characterized by a definitional band and
  a negligible estimation coupling, not a certified value.
\end{itemize}

None of these caveats touches the proven content---the classification, its basis invariance,
the exact defect count, the necessary-not-sufficient criterion, and the lift theorem---or
the validated content---the linear-algebra reproduction on real DE440-orientation rows. Each
marks where a real lunar deployment must supply flight data and a fixed convention before any
absolute number is trusted.

\section{Conclusion}
\label{sec:conclusion}

A lunar PNT realization fixes a spatial frame and a timescale jointly, and this paper
classifies what such a realization fails to determine about the two together. Over a
nine-parameter coupled datum---translation, scale, rotation, a timescale offset $\dtau$, and
a timescale rate $\dalpha$---we prove a \emph{basis-invariant} classification of the coupled
Fisher null space into spatial, temporal, and genuinely coupled parts, expressed through the
intrinsic projector $\Pnull=\Umat\Umat^\top$ and its off-diagonal block $\Pst$. Under the SI
second and the $c$-defined metre---adopted up front to quarantine the definitional
unit-of-time/scale coupling and keep scale observable---the exact datum gauge of a fully
internal \emph{free-network} one-way lunar problem (mutual common-mode clocks, no external
tie) is \emph{eight}: three translations, three rotations, the offset, and the rate;
anchoring to known orbiter positions and a time reference reduces this to defect zero, the
validated case. Scale is observable, so we never treat it as a free datum direction. A
coupled defect is certified only by $\Pst\neq\bm0$; a non-block-diagonal Fisher is necessary
but not sufficient. External ties lift the temporal defects---a time tie removes $\dtau$, and
the relativistic rate tie removes $\dalpha$ given an independent scale tie---and, once the
rate tie is admitted, the joint gauge acquires a spatial--temporal direction and is
\emph{not} the direct sum of the spatial and temporal gauges.

We quantify the two frame--time couplings and state one of them as an explicit correction.
Coupling~1 (range$\leftrightarrow$clock-offset, unit tie
$1\,\text{m}\leftrightarrow3.336\,\text{ns}=1/c$) is operationally strong but a near-defect,
and at the datum level, in a geometrically diverse one-way network, two-way ranging provides
\emph{no} selective scale$\leftrightarrow$offset decoupling beyond data volume (lift ratio
$\approx1.00$, computed on the modelled network)---the strong degeneracy is a single-user,
poor-geometry effect that does not aggregate. Coupling~2 (relativistic rate$\leftrightarrow$frame) is fundamentally weak: the
frame estimation datum perturbs the lunar timescale rate by $\lesssim\!\sim\!2\times10^{-17}$,
six orders of magnitude below the $\pm1.7\times10^{-11}$ reference-surface definitional band;
the rate is gauge-robust to estimation and gauge-sensitive to definition.

Exactly one claim is externally \Validated{}: the coupled-Fisher linear-algebra pipeline
reproduces an independent numpy/LAPACK computation, to relative and absolute error below
$10^{-3}$, on rows carrying real DE440 lunar-orientation partials, for a full-rank
(defect-$0$) and an exact-rate-defect network. The geometry is \Modelled{}, the coupling
magnitudes are \Modelled{}, and the nonzero-coupling branch of the classifier is \Proven{} by
construction and unit-tested, not externally validated---stated precisely.

Two design laws follow: geometric diversity separates the frame scale from the timescale
offset; and the rate's residual gauge is definitional, not estimational. Together they give
the compact requirement (Theorem~\ref{thm:system}): \emph{a lunar reference-system standard
must fix both a spatial datum and a temporal convention}---neither geometry, the relativistic
model, nor any volume of data can supply what the other must. The proven and cited content
carries over unchanged to any real deployment; before any absolute rate or frame number is
trusted, a real lunar realization must supply flight data and a fixed convention. Every number
here is reproducible from the open engine by a single command.


\section*{Data and code availability}
\label{sec:availability}
The reference substrate for this paper is part of the open-source \code{kshana} engine
\citep{baweja2026kshana} (\url{https://github.com/ashfordeOU/kshana}; AGPL-3.0-only). The
nine-vector coupled datum, the four observation-row builders
\eqref{eq:oneway}--\eqref{eq:ratetie}, the preconditioned Fisher assembly
\code{assemble\_coupled\_info} \eqref{eq:fisher}, the basis-invariant classifier
\code{classify\_null\_space} (Definition~\ref{def:defect},
Proposition~\ref{prop:invariance}), the $\{s,\dtau\}$ Schur marginal
\code{coupled\_marginal\_eigs} (Section~\ref{sec:coupling1}), and the relativistic
Jacobian \code{rate\_frame\_jacobian} \eqref{eq:dads}--\eqref{eq:dadr} are in the module
\code{lunar\_gauge}. Every number in this paper is regenerated by
\begin{quote}
  \code{cargo run -{}-example p4\_coupled\_gauge}
\end{quote}
which prints Part~A (the basis-invariant classification of a direct-sum and a coupled
matrix, plus the invariance demonstration), Part~B (the two real DE440-orientation anchor
networks of Table~\ref{tab:anchor} and the two-way-vs-more-one-way lift ratio of
Table~\ref{tab:coupling1}), and Part~C (the frame$\to$rate Jacobian, the frame-estimation
bound \eqref{eq:framebound}, the reference-surface definitional band \eqref{eq:band}, and
the $1\,\text{m}\leftrightarrow3.336\,\text{ns}$ unit tie), directly from the committed
engine and the vendored fixtures. No result is hard-coded in a print statement.

The single \Validated{} row---the coupled-Fisher pipeline matching an independent
numpy/LAPACK oracle to relative and absolute error below $10^{-3}$ on real
DE440-orientation rows---is asserted in the continuous-integration suite by
\begin{quote}
  \code{cargo test -{}-test lunar\_coupled\_gauge\_reference}
\end{quote}
which fails the build on any mismatch. The fixtures are regenerated by
\code{cargo run -{}-example gen\_coupled\_gauge\_rows} (which builds the observation rows
through the real DE440 lunar physical-axis orientation path and rounds every float to six
decimals) followed by the numpy oracle \code{scripts/gen\_coupled\_gauge\_ref.py}; the
vendored artifacts are \code{tests/fixtures/coupled\_gauge/} (\code{network.json},
\code{reference.json}, \code{NOTICE.md}), whose notice records the exact scope of what is
validated and what is modelled. The DE440 lunar orientation is JPL's
\citep{park2021de440}; the kernels are not redistributed. A productized operational lunar
timescale/frame realization---as opposed to the open reference substrate that reproduces the
paper's numbers---is a commercial capability of \code{kshana-pro} and is not part of this
open release.


\section*{Acknowledgements}
The numerical substrate underlying this work is part of the open-source
\code{kshana} engine. The lunar orientation and ephemeris data used here are the
JPL DE440 development ephemeris; the relativistic clock-rate coefficients are
derived from first principles and cross-checked against the published lunar
timescale band. This research received no specific grant from any funding agency
in the public, commercial, or not-for-profit sectors.

\bibliographystyle{plainnat}
\bibliography{refs}

\appendix
\section{Methods}
\label{app:methods}

This appendix records the numerical methods behind the results: the Fisher assembly and its
preconditioning, the Gram-rank threshold used by the classifier, the projector-based
classification algorithm, the $\{s,\dtau\}$ Schur complement, the independence of the
numpy/LAPACK oracle, and the derivation of the relativistic frame$\to$rate Jacobian. All are
consistent with the vendored provenance notice (Section~\ref{sec:availability}).

\subsection{Fisher assembly and preconditioning}
\label{app:fisher}

Each observation block $(\{\bm g_i\},\sigma)$ contributes
$\sigma^{-2}\sum_i\bm g_i\bm g_i^\top$ to the $9\times9$ Fisher \eqref{eq:fisher}, assembled
by \code{assemble\_coupled\_info}. Before accumulation, each row is preconditioned by
dividing columns $3$--$6$ (scale and the three rotations) by $\rmoon=\SI{1737400}{m}$,
bringing all nine partials to $\mathcal{O}(1)$. This is a positive diagonal column scaling
$\Imat\mapsto\bm D\Imat\bm D$; by Lemma~\ref{lem:precond} it preserves the null-space
dimension and the classification, changing only the eigenvalue \emph{magnitudes} (and hence
the condition number). All eigenvalues, the null space, and the marginal reported in the
tables are computed on the preconditioned matrix, and by Lemma~\ref{lem:precond} equal those
of the physical information.

\subsection{The Gram-rank threshold}
\label{app:gram}

The classifier needs the ranks of the row-blocks $\Umat_{\Ssp}$ and $\Umat_{\Tsp}$ of the
orthonormal null basis (Definition~\ref{def:defect}). It computes each as the number of
eigenvalues of the Gram matrix $\bm A\bm A^\top$ exceeding an absolute threshold $10^{-9}$,
where $\bm A$ is the block. Because $\Umat$ has orthonormal \emph{columns}, the Gram
eigenvalues of any row-block lie in $[0,1]$, so $10^{-9}$ cleanly separates zero from
$\mathcal{O}(1)$ and the threshold is dimensionless and invariant under positive column
rescaling. The overall defect $d$ is
taken from \code{fim::crlb} with a \emph{relative} tolerance $\text{rel\_tol}=10^{-9}$: an
eigenvalue of the full Fisher below $\text{rel\_tol}\cdot\lambda_{\max}$ is treated as zero.
For the \code{single\_epoch} network this gives a threshold
$10^{-9}\times44.97\approx4.5\times10^{-8}$, six orders below the smallest nonzero eigenvalue
$0.1605$---a clean gap (Section~\ref{sec:networks}). A genuine near-defect with a
sub-threshold temporal projection would sit on the spatial/coupled boundary; the anchored
networks do not, so the counts are unambiguous.

\subsection{The projector-based classification algorithm}
\label{app:classify}

Given the Fisher, \code{classify\_null\_space} (i)~calls \code{fim::crlb} to obtain the
orthonormal null basis $\Umat$ (layout $\Umat[\text{row}][\text{col}]$) and the defect $d$,
returning the all-zero class if $d=0$; (ii)~forms $\Umat_{\Ssp}=\Umat[0{:}7,:]$ and
$\Umat_{\Tsp}=\Umat[7{:}9,:]$ and computes $r_{\Ssp}=\rank(\Umat_{\Ssp})$,
$r_{\Tsp}=\rank(\Umat_{\Tsp})$ by the Gram threshold of Appendix~\ref{app:gram};
(iii)~derives $\dim(\Nsp\cap\Ssp)=d-r_{\Tsp}$, $\dim(\Nsp\cap\Tsp)=d-r_{\Ssp}$, and
$\operatorname{coupled\,dim}=r_{\Ssp}+r_{\Tsp}-d$ (a debug assertion guards
$r_{\Ssp}+r_{\Tsp}\ge d$ against underflow); and (iv)~forms the off-diagonal block of the
projector $\Pnull=\Umat\Umat^\top$ restricted to rows $0$--$6$ and columns $7$--$8$,
$\Pst[i][j]=\sum_{c}\Umat[i][c]\,\Umat[j][c]$, and returns
$\lVert\Pst\rVert_{\mathrm F}$. By Proposition~\ref{prop:invariance} every output depends only
on $\Nsp$, not on the eigensolver's basis choice; the unit test
\code{classify\_is\_invariant\_under\_null\_basis\_rotation} rotates $\Umat$ by
$\theta=0.6$~rad and confirms an identical class.

\subsection{The \texorpdfstring{$\{s,\dtau\}$}{\{s,dtau\}} Schur complement}
\label{app:schur}

The coupling-1 marginal (Section~\ref{sec:coupling1}) eliminates the seven parameters
$M=\{0,1,2,4,5,6,8\}$ (translations, rotations, and the rate) to expose the $2\times2$
information on $K=\{3,7\}=\{s,\dtau\}$. With the block partition of the Fisher into
$\Imat_{KK},\Imat_{KM},\Imat_{MM}$, the Schur complement is
\begin{equation}
  \bm S \;=\; \Imat_{KK}-\Imat_{KM}\,\Imat_{MM}^{+}\,\Imat_{KM}^\top,
  \label{eq:schur}
\end{equation}
where $\Imat_{MM}^{+}$ is the Moore--Penrose pseudo-inverse, taken from the
\code{fim::crlb} pseudo-covariance field, which remains finite even when the
$M$-block is rank-deficient (sparse networks). The complement \eqref{eq:schur} is a genuine
marginal information only when $\Imat_{MM}$ is full-rank; for a rank-deficient $\Imat_{MM}$
the pseudo-inverse makes $\bm S$ a \emph{regularized surrogate} that need not be PSD (a
negative marginal eigenvalue in an early single-beacon fixture was exactly this artifact).
All networks reported here have full-rank, positive-definite $M$-blocks, so $\bm S$ is a true
marginal and no reported value is affected. The eigenvalues of $\bm S$ (via
\code{fim::sym\_eig}, ascending) are reported as \code{marginal\_eigs}; $\lambda_{\min}(\bm S)$
is the coupling-1 near-defect strength. For the \code{well\_posed} network,
$\bm S$ is positive-definite with eigenvalues $[1.8720,\,35.9995]$.

\subsection{Independence of the numpy/LAPACK oracle}
\label{app:oracle}

The oracle \code{scripts/gen\_coupled\_gauge\_ref.py} reproduces the pipeline through a
\emph{different} numerical stack: eigenvalues by \code{np.linalg.eigvalsh}, null-space Gram
ranks by \code{np.linalg.eigh}, and the Schur pseudo-inverse by \code{np.linalg.pinv} (SVD),
against the engine's Jacobi eigensolver (\code{fim::sym\_eig}) and spectral-sum pseudo-inverse
(\code{fim::crlb}). The two backends---Jacobi versus LAPACK divide-and-conquer, spectral-sum
versus SVD pseudo-inverse---are independent implementations, so agreement to relative and
absolute error below $10^{-3}$ is a genuine cross-check, not a tautology. Non-circularity is
structural: the oracle reads the observation \emph{rows} from \code{network.json} and derives
\code{reference.json} from them, never from a Rust output; six-decimal rounding of the rows
before both writing and reading guarantees byte-identical inputs.

\subsection{Derivation of the frame\texorpdfstring{$\to$}{->}rate Jacobian}
\label{app:jacobian}

The three coefficients \eqref{eq:dads}--\eqref{eq:dadr} follow from the post-Newtonian clock
rate. We adopt the convention---used once here and throughout---that $\alpha$ is the rate of
a lunar-surface clock measured \emph{relative to Terrestrial Time} (TT), with TT tied to the
defining Earth geopotential $W_{0,\oplus}$. To first order a clock's rate is set by the
gravitational redshift and the second-order Doppler (kinetic) term, so
\begin{equation}
  \alpha \;\simeq\; \frac{W_{0,\oplus}-\Umoon}{c^2}-\frac{v^2}{2c^2}+\text{const},
  \label{eq:rate}
\end{equation}
with $\Umoon>0$ the lunar self-potential and $v$ the speed
\citep{soffel2003iau,mueller2018clocks}. In this surface-relative-to-TT convention the
self-potential enters with a \emph{minus} sign---a deeper lunar well runs the clock slower
relative to TT---so the coefficients below carry the $+\Umoon/c^2$ and $+\gmoon/c^2$ signs of
the $(W_0-U)/c^2$ form, opposite the sign of the bare redshift $\Phi=-GM/r$. For a clock on
the lunar surface the self-potential is $\Umoon=GM_{\mathrm{M}}/\rmoon$ and the surface
gravity is $\gmoon=GM_{\mathrm{M}}/\rmoon^{2}$. Differentiating:
\begin{itemize}
  \item \textbf{Scale.} A fractional scale change $\delta s$ scales the lunar radius,
  $\rmoon\mapsto\rmoon(1+\delta s)$, so the surface self-potential shrinks,
  $\Umoon(s)=\Umoon(1-s)$ to first order. Substituting into
  $\alpha=(W_{0,\oplus}-\Umoon(s))/c^2-v^2/(2c^2)+\text{const}$ gives
  $\partial\alpha/\partial s=+\Umoon/c^2=+3.140\times10^{-11}$ per unit fractional scale, with
  $GM_{\mathrm{M}}=4.9028\times10^{12}\,\mathrm{m^3/s^2}$ and
  $\rmoon=1.7374\times10^{6}$~m.
  \item \textbf{Radial.} A radial displacement $\delta r$ changes the potential by
  $\gmoon\,\delta r$, so $\partial\alpha/\partial r=+\gmoon/c^2=+1.807\times10^{-17}$ per
  metre. The identity $\partial\alpha/\partial r=(\partial\alpha/\partial s)/\rmoon$ is exact
  ($\gmoon/c^2=(\Umoon/\rmoon)/c^2$) and unit-tested to $10^{-12}$.
  \item \textbf{Velocity.} The kinetic term gives $\partial\alpha/\partial v=-|v|/c^2$, with
  $|v|$ the geocentric Moon speed at the epoch (from \code{lunar\_time}); at 2024-01-01
  ($|v|\approx969$~m/s) this is $-1.078\times10^{-14}$, near J2000 ($|v|\approx997$~m/s)
  $\approx-1.11\times10^{-14}$ per (m/s). The velocity coefficient is epoch-dependent by the
  several-percent monthly variation of the Moon's speed (the quoted $969$ vs $997$~m/s
  already differ by nearly $3\%$).
\end{itemize}
The external (Earth) tidal potential is \emph{excluded} from $\Umoon$: by the equivalence
principle a freely-falling body does not sense the uniform part of the external potential, so
$GM_{\oplus}/d$ does not enter the clock rate; only the tidal \emph{gradient} does, at
$\sim2.7\times10^{-22}$ per metre, negligible. All three coefficients are \Modelled{}
(first-principles, cross-checked to the $[56,59]\,\mu$s/day lunar-clock-rate figures
% confirm the [56,59] $\mu$s/day figure against fienga2024moonlight at camera-ready
consistent with the Moonlight reference-system literature \citep{fienga2024moonlight}), not
certified operational values.


\end{document}
